%% secoes/resultados/03-malformacao.tex — incluído por main.tex

%%---------------------------------------------------------------------%%
\subsection{Malformação congênita}\label{subsec:res-malf}
%%---------------------------------------------------------------------%%

A Tabela~\ref{tab:malf-dist} mostra a distribuição da malformação congênita nas oito populações. No recorte geral, a prevalência basal do desfecho foi de 0,60\%, a menor entre os quatro desfechos analisados.

As prevalências variam pouco entre os segmentos, de 0,47\% a 0,63\%. Seis dos sete segmentos vulneráveis ficaram abaixo do recorte geral: quilombola (0,47\%), indígena (0,48\%), rural (0,48\%), municípios de até 10 mil habitantes (0,54\%), mães com baixa escolaridade (0,56\%) e beneficiárias do Bolsa Família (0,56\%). Apenas as mães menores de 15 anos ficaram acima (0,63\%). Por se tratar de um desfecho raro, os segmentos menores reúnem poucas ocorrências: 227 na população indígena, 169 entre mães menores de 15 anos e 121 na população quilombola.

\begin{table}[htbp]
\caption{Distribuição da malformação congênita por população analisada.}\label{tab:malf-dist}%
\begin{tabularx}{\textwidth}{@{}Yrrrr@{}}
\toprule
População & Total & \makecell{Total\\(\%)} & \makecell{Nº de\\ocorrências} & \makecell{Ocorrências\\(\%)}\\
\midrule
Recorte geral & 5.246.679 & 100,00 & 31.270 & 0,60\\
Municípios até 10 mil hab. & 427.806 & 8,15 & 2.308 & 0,54\\
Rural & 1.133.241 & 21,60 & 5.446 & 0,48\\
Indígena & 47.474 & 0,90 & 227 & 0,48\\
Quilombola & 25.643 & 0,49 & 121 & 0,47\\
Mães com baixa escolaridade & 210.086 & 4,00 & 1.169 & 0,56\\
Mães menores de 15 anos & 26.671 & 0,51 & 169 & 0,63\\
Beneficiárias do Bolsa Família & 1.222.696 & 23,30 & 6.817 & 0,56\\
\bottomrule
\end{tabularx}
\fontetabela
\end{table}


A Tabela~\ref{tab:malf} apresenta os cinco subgrupos de maior WRAcc para a malformação congênita para os oito recortes considerados no artigo.
Na malformação congênita, dois itens aparecem nos oito recortes. O sexo masculino figura em pelo menos um subgrupo de cada recorte, e o parto cesáreo também. Em seis dos oito recortes, o subgrupo de maior risco contém o parto cesáreo: municípios de até 10 mil habitantes (31\% acima da basal), rural (47\%), indígena (71\%), quilombola (158\%), mães com baixa escolaridade (50\%) e mães menores de 15 anos (154\%). As exceções são o recorte geral e as beneficiárias do Bolsa Família, cujos subgrupos de maior risco combinam o sexo masculino a outros itens.

\begin{table}[p]
\centering
\caption{Subgrupos para malformação congênita (prevalência basal: 0,60\%).}
\label{tab:malf}

\scriptsize
\setlength{\tabcolsep}{4pt}
\renewcommand{\arraystretch}{1.04} % Levemente mais compacto para compensar subgrupos com 4-5 condições

\begin{tabularx}{\textwidth}{@{}c Y r r r r r@{}}
\toprule
\textbf{\#} & \textbf{Subgrupo} & \textbf{Total} & \textbf{\makecell[r]{Total\\(\%)}} & \textbf{\makecell[r]{Nº de\\ocorr.}} & \textbf{\makecell[r]{Ocorr.\\(\%)}} & \textbf{Lift} \\
\midrule
\secaotabela{Recorte geral}
1 & água canalizada: sim; banheiro: sim; sexo: M & 2.049.407 & 39,1 & 15.020 & 48,0 & 1,23 \\
2 & banheiro: sim; escoamento sanitário: rede coletora & 2.216.543 & 42,2 & 15.930 & 50,9 & 1,21 \\
3 & água canalizada: sim; banheiro: sim; parto cesárea & 1.947.713 & 37,1 & 14.012 & 44,8 & 1,21 \\
4 & escoamento sanitário: rede coletora; sexo: M; transferência de renda: não; Bolsa Família: não & 792.340 & 15,1 & 6.768 & 21,6 & 1,43 \\
5 & sexo: M; parto cesárea; Bolsa Família: não & 976.618 & 18,6 & 7.777 & 24,9 & 1,34 \\

\secaotabela{Municípios de até 10 mil habitantes}
6 & banheiro: sim; despesa com transporte: 0; nascimento hospitalar; parto cesárea; sabe ler e escrever: sim & 142.432 & 33,3 & 939 & 40,7 & 1,22 \\
7 & banheiro: sim; sexo: M; despesa com transporte: 0; nascimento hospitalar & 148.130 & 34,6 & 953 & 41,3 & 1,19 \\
8 & sexo: M; parto cesárea; Bolsa Família: não & 81.234 & 19,0 & 570 & 24,7 & 1,30 \\
9 & despesa com alimentação: 100--300; parto cesárea; sabe ler e escrever: sim & 76.541 & 17,9 & 520 & 22,5 & 1,26 \\
10 & nascimento hospitalar; parto cesárea; primeira gestação: sim & 63.649 & 14,9 & 449 & 19,5 & 1,31 \\

\secaotabela{População rural}
11 & nascimento hospitalar; parto cesárea & 434.046 & 38,3 & 2.569 & 47,2 & 1,23 \\
12 & nascimento hospitalar; sexo: M & 552.742 & 48,8 & 3.082 & 56,6 & 1,16 \\
13 & banheiro: sim; sexo: M; parto cesárea & 165.018 & 14,6 & 1.164 & 21,4 & 1,47 \\
14 & banheiro: sim; iluminação: eletricidade & 680.578 & 60,1 & 3.553 & 65,2 & 1,09 \\
15 & nascimento hospitalar; primeira gestação: sim & 270.911 & 23,9 & 1.576 & 28,9 & 1,21 \\

\secaotabela{População indígena}
16 & nascimento hospitalar; parto cesárea; sabe ler e escrever: sim & 7.943 & 16,7 & 65 & 28,6 & 1,71 \\
17 & abastecimento de água: poço/nascente; nascimento hospitalar & 14.657 & 30,9 & 96 & 42,3 & 1,37 \\
18 & sexo: M; despesa com transporte: 0; sabe ler e escrever: sim & 14.719 & 31,0 & 93 & 41,0 & 1,32 \\
19 & abastecimento de água: poço/nascente; sexo: M; despesa com transporte: 0 & 6.909 & 14,6 & 54 & 23,8 & 1,63 \\
20 & nascimento hospitalar; sabe ler e escrever: sim & 31.353 & 66,0 & 168 & 74,0 & 1,12 \\

\secaotabela{População quilombola}
21 & material do domicílio: tijolo/alvenaria; material do piso: cerâmica/pedra/carpete; famílias no domicílio: 1 ou menos; raça/cor: parda/preta & 9.239 & 36,0 & 64 & 52,9 & 1,47 \\
22 & banheiro: sim; material do piso: cerâmica/pedra/carpete; dormitórios: 2--3; famílias no domicílio: 1 ou menos; despesa com água/esgoto: 0 & 4.942 & 19,3 & 40 & 33,1 & 1,72 \\
23 & banheiro: sim; material do piso: cerâmica/pedra/carpete; dormitórios: 2--3; famílias no domicílio: 1 ou menos; parto cesárea & 2.218 & 8,6 & 27 & 22,3 & 2,58 \\
24 & dormitórios: 2--3; famílias no domicílio: 1 ou menos; despesa com energia: 0--20 & 5.220 & 20,4 & 39 & 32,2 & 1,58 \\
25 & material do piso: cerâmica/pedra/carpete; dormitórios: 2--3; famílias no domicílio: 1 ou menos; sexo: M & 5.462 & 21,3 & 40 & 33,1 & 1,55 \\

\secaotabela{Mães com baixa escolaridade}
26 & material do domicílio: tijolo/alvenaria; sexo: M & 71.325 & 34,0 & 525 & 44,9 & 1,32 \\
27 & escoamento sanitário: rede coletora; iluminação: eletricidade & 54.802 & 26,1 & 389 & 33,3 & 1,28 \\
28 & material do domicílio: tijolo/alvenaria; sexo: M; parto cesárea & 28.447 & 13,5 & 238 & 20,4 & 1,50 \\
29 & material do domicílio: tijolo/alvenaria; famílias no domicílio: 1 ou menos & 134.986 & 64,3 & 828 & 70,8 & 1,10 \\
30 & parto cesárea; famílias no domicílio: 1 ou menos & 70.558 & 33,6 & 457 & 39,1 & 1,16 \\

\secaotabela{Mães menores de 15 anos}
31 & abastecimento de água: rede pública; água canalizada: sim; cômodos no domicílio: 3--4 & 7.762 & 29,1 & 76 & 45,0 & 1,55 \\
32 & material do domicílio: tijolo/alvenaria; sexo: M; primeira gestação: sim & 7.329 & 27,5 & 66 & 39,1 & 1,42 \\
33 & água canalizada: sim; cômodos no domicílio: 3--4; parto cesárea & 3.714 & 13,9 & 43 & 25,4 & 1,83 \\
34 & abastecimento de água: rede pública; água canalizada: sim; material do domicílio: tijolo/alvenaria; raça/cor da mãe: branca & 4.158 & 15,6 & 45 & 26,6 & 1,71 \\
35 & escoamento sanitário: rede coletora; cômodos no domicílio: 3--4; parto cesárea & 1.803 & 6,8 & 29 & 17,2 & 2,54 \\

\secaotabela{Beneficiárias do Programa Bolsa Família}
36 & banheiro: sim; destino do lixo: coletado; domicílio particular permanente; sexo: M & 418.152 & 34,2 & 2.929 & 43,0 & 1,26 \\
37 & escoamento sanitário: rede coletora; domicílio particular permanente & 433.860 & 35,5 & 2.889 & 42,4 & 1,19 \\
38 & sexo: M; parto cesárea & 273.979 & 22,4 & 1.980 & 29,0 & 1,30 \\
39 & sexo: M; material do piso: madeira/cimento & 204.150 & 16,7 & 1.477 & 21,7 & 1,30 \\
40 & sexo: M; água canalizada: sim; dormitórios: 1 ou menos & 172.561 & 14,1 & 1.299 & 19,1 & 1,35 \\
\bottomrule
\end{tabularx}
%\fontetabela
\end{table}

%geral
%No recrte geral, o banheiro no domicílio e o sexo masculino aparecem em três subgrupos cada (1, 2 e 3; 1, 4 e 5). Água canalizada, escoamento por rede coletora, parto cesáreo e ausência de Bolsa Família aparecem em dois subgrupos cada.

%O subgrupo 4 apresenta o maior risco do desfecho da tabela (43\% acima da basal) e a menor participação nas ocorrências (21,6\%). Os subgrupos 2 e 3 têm o menor risco do desfecho (21\% acima da basal), e o subgrupo 2 tem a maior participação nas ocorrências (50,9\%).

%Os três primeiros subgrupos combinam itens de infraestrutura domiciliar presente (banheiro, água canalizada e escoamento por rede coletora) e abrangem, cada um, mais de 37\% dos nascimentos e mais de 44\% das ocorrências, com risco entre 21\% e 23\% acima da basal. Os subgrupos 4 e 5, que têm em comum o sexo masculino e a ausência de Bolsa Família, são menores e apresentam risco mais alto (43\% e 34\% acima da basal).

%10k
%No municípios de até 10 mil habitantes, o parto cesáreo aparece em quatro dos cinco subgrupos (6, 8, 9 e 10) e o nascimento hospitalar em três (6, 7 e 10). Banheiro no domicílio, despesa com transporte nula, mãe que sabe ler e escrever e sexo masculino aparecem em dois subgrupos cada.

%O subgrupo 10 (nascimento hospitalar, parto cesáreo e primeira gestação) apresenta o maior risco do desfecho da tabela (31\% acima da basal) e a menor participação nas ocorrências (19,5\%). O subgrupo 7 tem o menor risco do desfecho (19\% acima da basal) e a maior participação nas ocorrências (41,3\%). Os riscos variam pouco, entre 19\% e 31\% acima da basal.

%Rural
%Na população rural, o nascimento hospitalar aparece em três subgrupos (11, 12 e 15). Parto cesáreo, sexo masculino e banheiro no domicílio aparecem em dois subgrupos cada.

%O subgrupo 13 (banheiro, sexo masculino e parto cesáreo) apresenta o maior risco do desfecho da tabela (47\% acima da basal) e a menor participação nas ocorrências (21,4\%). O subgrupo 14 (banheiro e iluminação elétrica) tem o menor risco do desfecho (9\% acima da basal) e a maior participação nas ocorrências (65,2\%), abrangendo 60,1\% dos nascimentos do segmento.

%Os subgrupos 11, 12 e 15 compartilham o nascimento hospitalar e diferem no segundo item: com parto cesáreo, o risco é de 23\% acima da basal; com primeira gestação, de 21\%; e com sexo masculino, de 16\%.

%indigena
%Na população indígena, o nascimento hospitalar e a mãe que sabe ler e escrever aparecem em três subgrupos cada (16, 17 e 20; 16, 18 e 20). Abastecimento de água por poço ou nascente, sexo masculino e despesa com transporte nula aparecem em dois subgrupos cada.

%O subgrupo 16 (nascimento hospitalar, parto cesáreo e mãe que sabe ler e escrever) apresenta o maior risco do desfecho da tabela (71\% acima da basal). A menor participação nas ocorrências é a do subgrupo 19 (23,8\%). O subgrupo 20 tem o menor risco do desfecho (12\% acima da basal) e a maior participação nas ocorrências (74,0\%), abrangendo 66,0\% dos nascimentos do segmento.

%Os subgrupos 20 e 16 diferem por um único item: a adição do parto cesáreo eleva a prevalência de 0,54\% para 0,82\%, de 12\% para 71\% acima da basal, com queda da participação nas ocorrências de 74,0\% para 28,6\%.

%quilombola

%Na população quilombola, o item famílias no domicílio (1 ou menos) aparece nos cinco subgrupos, o material do piso (cerâmica, pedra ou carpete) e o número de dormitórios (2 a 3) em quatro, e o banheiro no domicílio em dois (2 e 3).

%O subgrupo 23 apresenta o maior risco do desfecho da tabela (158\% acima da basal, isto é, 2,58 vezes a basal), com prevalência de 1,22\% e a menor participação nas ocorrências (22,3\%), o que corresponde a 27 casos. O subgrupo 21 tem o menor risco do desfecho (47\% acima da basal) e a maior participação nas ocorrências (52,9\%).

%Os subgrupos 22 e 23 compartilham quatro itens (banheiro, material do piso, dormitórios e famílias no domicílio) e diferem no quinto: com despesa com água e esgoto nula, o risco é de 72\% acima da basal; com parto cesáreo, de 158\%.

%maes baixa escolaridade
%Nas mães com baixa escolaridade, o material do domicílio (tijolo ou alvenaria) aparece em três subgrupos (26, 28 e 29). Sexo masculino, parto cesáreo e famílias no domicílio (1 ou menos) aparecem em dois subgrupos cada.

%O subgrupo 28 apresenta o maior risco do desfecho da tabela (50\% acima da basal) e a menor participação nas ocorrências (20,4\%). O subgrupo 29 tem o menor risco do desfecho (10\% acima da basal) e a maior participação nas ocorrências (70,8\%), abrangendo 64,3\% dos nascimentos do segmento.

%Os subgrupos 26 e 28 diferem por um único item: a adição do parto cesáreo eleva a prevalência de 0,74\% para 0,84\%, de 32\% para 50\% acima da basal, com queda da participação nas ocorrências de 44,9\% para 20,4\%.

%maes menores
%Nas mães menores de 15 anos, a água canalizada e o número de cômodos (3 a 4) aparecem em três subgrupos cada (31, 33 e 34; 31, 33 e 35). Parto cesáreo, abastecimento de água por rede pública e material do domicílio (tijolo ou alvenaria) aparecem em dois subgrupos cada.

%O subgrupo 35 apresenta o maior risco do desfecho da tabela (154\% acima da basal, isto é, 2,54 vezes a basal), com prevalência de 1,61\% e a menor participação nas ocorrências (17,2\%), o que corresponde a 29 casos. O subgrupo 32 tem o menor risco do desfecho (42\% acima da basal), e o subgrupo 31 tem a maior participação nas ocorrências (45,0\%).

%Os subgrupos 33 e 35 compartilham o número de cômodos e o parto cesáreo e diferem no terceiro item: com água canalizada, o risco é de 83\% acima da basal; com escoamento por rede coletora, de 154\%.

%Nos beneficiárias do Bolsa Família, o sexo masculino aparece em quatro dos cinco subgrupos (36, 38, 39 e 40) e o domicílio particular permanente em dois (36 e 37).

%O subgrupo 40 (sexo masculino, água canalizada e até um dormitório) apresenta o maior risco do desfecho da tabela (35\% acima da basal) e a menor participação nas ocorrências (19,1\%). O subgrupo 37 tem o menor risco do desfecho (19\% acima da basal), e o subgrupo 36 tem a maior participação nas ocorrências (43,0\%). Os riscos variam pouco, entre 19\% e 35\% acima da basal.

%tudo 
%Na malformação congênita, dois itens aparecem nos oito recortes. O sexo masculino figura em pelo menos um subgrupo de cada recorte, e o parto cesáreo também. Em seis dos oito recortes, o subgrupo de maior risco contém o parto cesáreo: municípios de até 10 mil habitantes (31\% acima da basal), rural (47\%), indígena (71\%), quilombola (158\%), mães com baixa escolaridade (50\%) e mães menores de 15 anos (154\%). As exceções são o recorte geral e as beneficiárias do Bolsa Família, cujos subgrupos de maior risco combinam o sexo masculino a outros itens.

Os itens de infraestrutura domiciliar aparecem em todos os recortes e, em sua maioria, descrevem condições presentes: banheiro, água canalizada, escoamento por rede coletora, abastecimento por rede pública, iluminação elétrica e domicílio de tijolo ou alvenaria. A exceção é a população indígena, em que o abastecimento de água por poço ou nascente figura em dois subgrupos. O nascimento hospitalar aparece nos municípios de até 10 mil habitantes, no rural e no indígena.

O número de consultas de pré-natal não figura em nenhum dos 40 subgrupos do desfecho.
Algumas regras se repetem entre recortes com riscos próximos. A conjunção de sexo masculino, parto cesáreo e ausência de Bolsa Família aparece no recorte geral (34\% acima da basal) e nos municípios de até 10 mil habitantes (30\%); sexo masculino e parto cesáreo, sem o item do Bolsa Família, aparece entre as beneficiárias do programa (30\%). A conjunção de nascimento hospitalar e primeira gestação está 21\% acima da basal no rural, e a adição do parto cesáreo a esses dois itens, nos municípios de até 10 mil habitantes, corresponde a 31\%.

Os maiores riscos do desfecho estão em dois segmentos pequenos: o subgrupo 23 da população quilombola (158\% acima da basal) e o subgrupo 35 das mães menores de 15 anos (154\%), ambos com parto cesáreo e ambos com menos de 30 ocorrências (27 e 29). Nos recortes maiores, os riscos máximos são mais baixos: 31\% nos municípios de até 10 mil habitantes, 35\% entre beneficiárias do Bolsa Família e 43\% no recorte geral. Em três recortes (rural, indígena e mães com baixa escolaridade), um único subgrupo abrange mais de 60\% dos nascimentos, com risco entre 9\% e 12\% acima da basal.

\FloatBarrier
